% ID: SEL-D0035-OA-E03
% ESTADO: cotejada. Segunda pasada visual y borrador compilado revisados.
% Fuente: Selectividad/matematicas_aplicadas/titular_septiembre_EXAMEN 1-ANDALUCÍA.pdf
% SHA-256: 7a09a02641ea8de345b94f49b27515e58436e032c7401867f86dc523490a3ea5
% Página: 1; opción A; ejercicio original 3.

Supongamos que el 20\% de los votantes de Trump apoya la construcción del muro en la frontera con México y que solo el 5\% de los que no lo votaron la apoya. En un grupo formado por 5000 votantes de Trump y 10000 estadounidenses que no lo votaron se elige una persona al azar.
\begin{enumerate}[label=\alph*)]
\item \textbf{(1 punto)} \textquestiondown{}Cuál es la probabilidad de que ésta apoye la construcción del muro?
\item \textbf{(0.75 puntos)} Si la persona elegida apoya la construcción del muro, \textquestiondown{}cuál es la probabilidad de que no haya votado a Trump?
\item \textbf{(0.75 puntos)} Calcule la probabilidad de que sea votante de Trump o apoye la construcción del muro.
\end{enumerate}
